\documentclass[acmsmall, review, screen, authorversion]{acmart}

\usepackage{
    amsmath,mathpartir,mathtools,minted,newunicodechar,simplebnf,
    xcolor
}

\newunicodechar{∀}{$\forall$}
\newunicodechar{∷}{$\dblcolon$}

% \acmConference[ICFP'26]
%     {International Conference on Functional Programming}
%     {August 24--29, 2026}{Indianapolis, IN, USA}
% \acmBooktitle{Proceedings of the ACM on Programming Languages}

\title{Type Systems That Reinterpret Themselves}
\author{Pavel Sokolov}
\email{ppsokolov@hse.ru}

\begin{CCSXML}
<ccs2012>
   <concept>
       <concept_id>10003752.10003790.10011740</concept_id>
       <concept_desc>Theory of computation~Type theory</concept_desc>
       <concept_significance>500</concept_significance>
       </concept>
 </ccs2012>
\end{CCSXML}

\ccsdesc[500]{Theory of computation~Type theory}

\keywords{
    components and composition, type systems, dependent types,
    type theory, category theory
}

\begin{abstract}
    We advocate that type systems should be able to reinterpret their
    own terms in any internal model of themselves: equipped with a
    single tool of reinterpretation, you can express many desired
    features of type systems, including universe polymorphism,
    limited effect polymorphism, external parametricity and more.
    To this end, we present two type systems capable of
    self-reinterpretation, namely STLC-R and MLTT-R, prove their
    basic properties and show how the features we claim they contain
    are definable on the library level.
\end{abstract}

\begin{document}

\maketitle

\section{Introduction}

Consider the following well-known function which maps over elements
of the list, written in Agda:

\begin{minted}{agda}
map : ∀ {A B : Set} → (A → B) → List A → List B
map f [] = []
map f (x ∷ xs) = f x ∷ map f xs
\end{minted}

Enjoying both composability of this function and simplicity of its
type, we would like to reuse its definition as much as possible. Up
to a certain level, we would surely be able to do so: after all, its
definition is polymorphic in source and target types, so we can
instantiate it with whatever types we want. But we might wish for
much more; of things that come to our mind are the following:

\begin{itemize}
    \item Reuse \texttt{map} at higher universe levels, e.g., given a
        \textit{type operator}, map over \textit{list of types};
    \item Reuse \texttt{map} with a mapping function which is
        effectful, e.g. throws exceptions, logs something to the
        output etc.;
    \item Prove something about \texttt{map} for free, obtaining free
        theorem via parametricity translation.
\end{itemize}

The first point is actually achievable by adjusting the definition
and making it level-polymorphic. Second one, too, by generalizing
over side effects, probably specifying that it should work with any
\texttt{Applicative} functor, borrowing terminology from Haskell; or
maybe we could imagine that we have an effect system and specify that
function works with arbitrary effect $e$. Applying all of these
changes, we get something else entirely:

\begin{minted}{agda}
traverse : ∀ {a b} {F : Set b → Set b} {{Applicative F}}
         → ∀ {A : Set a} {B : Set b}
         → (A → F B) → List A → F (List B)
traverse f [] = pure []
traverse f (x ∷ xs) = liftA2 (∷) (f x) (traverse f xs)
\end{minted}

Note how the actual behaviour of the function, specified in a type,
gets obscured by the bookkeeping; getting a function to behave like
before is also harder as well, since its API became harder to use.
Not to mention that we didn't get a free theorem we wanted to. This
is possible in a system with \textit{internal parametricity}, as
exemplified by Altenkirch and others
\cite{Altenkirch:Parametricity2024}; yet, property extracted from
\texttt{traverse} would be different from the one you get from
\texttt{map}.

A shift in perspective might help. Instead of tweaking the definition
of \texttt{map}, let's remember all of the metatheoretic power our
type system already has. Parametricity is a first hint; another is
due to Lambek \cite{Lambek1980} who established that simply-typed
lambda calculus (STLC) is an \textit{internal language} of
\textit{cartesian closed categories} (CCC) \footnote{This is, of
course, part of well-known Curry-Howard-Lambek correspondence.}. In
particular, any term written in STLC can be interpreted in any
cartesian closed category; this way, you don't have to repeat your
construction in any CCC you meet, you just reinterpret the STLC term
you already have \cite{Elliott:Compcat2017}. This correspondence
extends further into dependently typed systems, establishing
correspondence with certain kinds of topoi and enabling us to do
synthetic mathematics: prove general statements in higher-order
intuitionistic logic and translate them into particular theorems
about their topos of interest \cite{Nawrocki:Lean2026}.

All of this power, unfortunately, lies beyond what usual type systems
can see, since that's their metatheoretical properties. But what if
they were internalized? What if we had a way to \textit{reinterpret
terms of a type system inside itself}? We experiment with this idea
and end up with two minimal systems, STLC-R (section \ref{STLC-R})
and MLTT-R (section \ref{MLTT-R}), based on simply typed lambda
calculus and Martin-L\"of type theory respectively. Turns out that
their capability of \textit{self-reinterpretation} is expressive
enough to cover all of the discussed usecases, including restricted
form of effect polymorphism (section \ref{STLC-R-applications}),
synthetic mathematics (section \ref{MLTT-R-synthetic}), universe
polymorphism (section \ref{MLTT-R-displacement}), parametricity
(section \ref{MLTT-R-parametricity}) and more.

\subsection{Contributions}

\begin{itemize}
    \item We start by giving the syntax and semantics of our first
        type system that reinterprets itself. It is based on simply
        typed lambda calculus and is called STLC-R. We prove that
        STLC-R is sound, and that global type inference is decidable
        (section \ref{STLC-R}).
    \item As expected, you can reinterpret STLC-R terms in any
        cartesian closed category you define inside it, including
        categories of differentiable manifolds for applications in
        autodifferentiation, Kleisli categories of certain monads and
        more (section \ref{STLC-R-applications}).
    \item Building upon the insights gathered from developing STLC-R,
        we give the syntax and semantics of a dependently typed
        self-reinterpreting language based on Martin-L\"of type
        theory called MLTT-R. We prove soundness of MLTT-R as well,
        and motivate bidirectional typing discipline for it (section
        \ref{MLTT-R}).
    \item As we intended, MLTT-R incorporates both a restricted form
        of universe polymorphism and external parametricity by virtue
        of reinterpretation only. And, of course, you can use MLTT-R
        as a system for synthetic mathematics with internal reuse of
        synthetic results (section \ref{MLTT-R-applications}).
\end{itemize}

\section{Almost simply typed reinterpretable language: STLC-R}
\label{STLC-R}

We start with a simpler language to iron out the details and actually
figure out how to properly reinterpret terms internally before moving
on to dependently typed system. Our starting point is simply typed
lambda calculus with products, which we call <<STLC>> throughout this
paper. As discussed before, we will add self-reinterpretation
capabilities to it to enable reinterpretation in any <<model>> of
STLC specifiable inside our language (here, we call them
<<universes>>). Here's the thing we \textit{wish} to do with it,
where $t\,@\,\mathcal{U}$ is an operator of reinterpretation of term
$t$ in internal model $\mathcal{U}$:

\begin{multline*}
    (\lambda f x y. (f x, y))\,@\,\langle
        (\circ),\mathtt{pair},\mathtt{proj}_1,\mathtt{proj}_2,
        \mathtt{lam},\mathtt{eval},\ldots
    \rangle=\\=
    \mathtt{lam}\,(\mathtt{lam}\,(\mathtt{lam}\,(\mathtt{pair}\,(
        \mathtt{eval}\,\circ\,\mathtt{pair}\,(
            \mathtt{proj}_2\circ\mathtt{proj}_1\circ\mathtt{proj}_1
        )\,(\mathtt{proj}_2\circ\mathtt{proj}_1)
    )\,\mathtt{proj}_2)))
\end{multline*}

Where $\mathtt{lam}$,\,$\mathtt{pair}$ and $\mathtt{eval}$ are
correspondingly constructor of an exponential object, constructor of
a product and evaluation morphism, in the same way they are specified
in the definition of cartesian closed categories.

There are a few things to ponder:

\begin{itemize}
    \item How to reinterpret free variables?
    \item How to type universes?
    \item We're in a simply typed setting, so there's no way to
        enforce that universe satisfies coherence conditions on their
        operations. How to structure computation in such a way that
        it doesn't break Church-Rosser property?
\end{itemize}

To answer the first question, we actually disallow (most of) free
variables in a terms that are to be reinterpreted since it is not
clear how to lift arbitrary \textit{value} into current universe, so
we have to enforce some kind of a typing discipline for introducing
terms amenable to reinterpretation. In a sense, reinterpretation is a
strong form of polymorphism, so let's structure our grammar and
typing rules similarly to how it's done in Hindley-Milner type
systems today and see how far we can go.

\begin{figure}[h]
    \centering
    \begin{bnf}
        $S,T$ : \textsf{MonoType} ::=
        | $\alpha$ : type variable
        | $\top$ : unit type
        | $S\times T$ : product type
        | $S\to T$ : arrow type
        ;;
        $\sigma$ : \textsf{PolyType} ::=
        | T : monomorphic type
        | $\forall\alpha.\sigma$ : forall type 
        ;;
        $s,t$ : \textsf{Term} ::=
        | $x$ : variable
        | $()$ : singleton
        | $(s, t)$ : pair
        | $\pi_1\,t$ : first projection
        | $\pi_2\,t$ : second projection
        | $\lambda x. t$ : abstraction
        | $t\,s$ : application
        | \texttt{let}\,$x=s$\,\texttt{in}\,$t$ : polymorphic binding
        | $x\,@\,\mathcal{U}^\cdot$ : reinterpretation
        ;;
        $\mathcal{U}$ : \textsf{Universe} ::=
        | $\langle T, t, \ldots \rangle$ : universe def. (see below)
    \end{bnf}
    \caption{STLC-R grammar}
    \Description{STLC-R grammar}
    \label{stlc-r-grammar}
\end{figure}

\begin{figure}[h]
    \centering
    \begin{mathpar}
        \inferrule{ }{\Gamma\vdash ():\top}
        \and\inferrule
            {\Gamma\vdash s:S\\\Gamma\vdash t:T}
            {\Gamma\vdash (s,t):S\times T}
        \and\inferrule
            {\Gamma\vdash t:S\times T}
            {\Gamma\vdash \pi_1\,t:S}
        \and\inferrule
            {\Gamma\vdash t:S\times T}
            {\Gamma\vdash \pi_2\,t:T}\\
        \and\inferrule{x:\sigma\in\Gamma}{\Gamma\vdash x:\sigma}
        \and\inferrule
            {\Gamma,x:S\vdash t:T}
            {\Gamma\vdash\lambda x.t:S\to T}
        \and\inferrule{
            \Gamma\vdash t:S\to T\\
            \Gamma\vdash s:S
        }{\Gamma\vdash t\,s:T}
        \and\inferrule{
            \forall\Gamma\vdash s:S\\
            \vec{\alpha}=fv(S)\\
            \forall\Gamma,x:\forall\vec{\alpha}.S,\Delta\vdash t:T
        }{\forall\Gamma,\Delta\vdash
        \mathtt{let}\,x=s\,\mathtt{in}\,t:T}
        \and\inferrule{
            \forall\Gamma\vdash x:\forall\vec{\alpha}.T\\
            \mathcal{U}\,\mathrm{universe}
        }{\forall\Gamma,\Delta\vdash x\mathtt{@}\mathcal{U}^\cdot
        :\mathcal{U}\llbracket T\rrbracket\big[
            \vec{\alpha}\coloneqq\vec{S}\big]}
        \and\inferrule{
            0,1\,\mathrm{MonoType}\vdash C,P,E\,\mathrm{MonoType}\\
            \mathtt{id}:C[\alpha,\alpha]\\
            \mathtt{comp}:
                C[\beta,\gamma]\to C[\alpha,\beta]
                \to C[\alpha,\gamma]\\
            TT\,\mathrm{MonoType}\\
            \mathtt{cnst}:C[\alpha,TT]\\
            \mathtt{pair}:
                C[\gamma,\alpha]\to C[\gamma,\beta]
                \to C[\gamma,P[\alpha,\beta]]\\
            \mathtt{proj}_1: C[P[\alpha,\beta],\alpha]\\
            \mathtt{proj}_2: C[P[\alpha,\beta],\beta]\\
            \mathtt{lam}:
                C[P[\alpha,\beta],\gamma]
                \to C[\alpha,E[\gamma,\beta]]\\
            \mathtt{eval}:
                C[P[E[\alpha,\beta],\beta],\alpha]
        }{\langle
            \Lambda 01.C,\mathtt{id},\mathtt{comp},TT,\mathtt{cnst},
            \Lambda 01.P,\mathtt{pair},\mathtt{proj}_1,
                \mathtt{proj}_2,
            \Lambda 01.E,\mathtt{lam},\mathtt{eval}
        \rangle\,\mathrm{universe}}
    \end{mathpar}
    \caption{STLC-R declarative typing rules}
    \Description{STLC-R declarative typing rules}
    \label{stlc-r-typing}
\end{figure}

In figure \ref{stlc-r-typing}, we mostly repeat usual declarative
syntax-directed typing rules of a Hindley-Milner type system
\cite{Clement:ML1986}, with important differences in \texttt{let}
case and in the case of reinterpretation:

\begin{itemize}
    \item As discussed, we don't take local context in account when
        checking polymorphic term $s$. Instead, it can only reference
        other polymorphic terms. This is safe since other polymorphic
        terms are constrained in the same way, and can be safely
        reinterpreted inside $s$.
    \item Polymorphic terms have to be instantiated with a universe
        when used, so they get a separate rule instead of usual rule
        for variables, which is only used for monomorphic variables.
\end{itemize}

To check universes, we abuse our type system a bit in order to avoid
adding type operators and complicate the type system: types of
arrows, products and exponentials are specified as open monomorphic
types with two type variables $0$,$1$ in a context which are then
filled in by substitution. Operations in a universe are typed with
types having type variables $\alpha$,$\beta$,$\gamma$ in a context
once again, since we don't have ``usual'' polymorphic functions, only
reinterpreting ones.

We could generalize universe definition to include something else as
objects instead of types, but there is no way to build meaningful
$C$, $P$ and $E$ with anything other than types in base system.

\begin{definition}
    Let $\mathcal{U}=\langle
        \Lambda 01.C,\mathtt{id},\mathtt{comp},TT,\mathtt{cnst},
        \Lambda 01.P,\mathtt{pair},\mathtt{proj}_1,\mathtt{proj}_2,
        \Lambda 01.E,\mathtt{lam},\mathtt{eval}
    \rangle$. Then, interpretation of type in a universe is the
    following function:
    \begin{align*}
        \mathcal{U}\llbracket T\rrbracket
        &= C[TT,T_\mathcal{U}]\\
        \top_\mathcal{U} &= TT\\
        (S\times T)_\mathcal{U}
        &= P[S_\mathcal{U},T_\mathcal{U}]\\
        (S\to T)_\mathcal{U}
        &= E[S_\mathcal{U},T_\mathcal{U}].
    \end{align*}
\end{definition}

Having defined typing rules, let's pin down computation rules in
order to prove some useful properties. General idea is the following:
inside \texttt{let} bindings, we want to only reduce other
$\texttt{let}$ bindings; in all other cases, standard reductions of
STLC apply. Proper definition is given a bit later, here we list only
the most meaningful rules.

\begin{figure}[h]
    \centering
    \begin{mathpar}
    \inferrule{}{\pi_i(t_1, t_2)\to_\beta t_i}
    \and\inferrule{}{(\lambda x.t)s\to_\beta t[x\coloneqq s]}\\
    \and\inferrule{s\to_\mathtt{let} s'}{
        \mathtt{let}\,x=s\,\mathtt{in}\,t
        \to_\mathtt{let}
        \mathtt{let}\,x=s'\,\mathtt{in}\,t
    }
    \and\inferrule{
        \mathtt{let}\not\in s
    }{
        \mathtt{let}\,x=s\,\mathtt{in}\,t
        \to_\mathtt{let} t[x\coloneqq s]
    }
    \end{mathpar}
    \caption{STLC-R computation rules}
    \Description{STLC-R computation rules}
    \label{stlc-r-comp}
\end{figure}

In computation rules given in fig.\ref{stlc-r-comp}, attentive reader
might notice that we might obtain invalid term after substitution
performed during \texttt{let}-reduction since only variables can be
reinterpreted in our grammar. Actually, by substitution into
$x\,@\,\mathcal{U}$ we mean the following:

\begin{definition}
\begin{align*}
(x\,@\,\mathcal{U})[x\coloneqq s] &= s\,@\,\mathcal{U}\\
()@\mathcal{U}^\gamma &= \mathtt{cnst}\\
(s,t)@\mathcal{U}^\gamma &= \mathtt{pair}\,
    (s@\mathcal{U}^\gamma)\,(t@\mathcal{U}^\gamma)\\
(\pi_i\,t)@\mathcal{U}^\gamma &= \mathtt{comp}\,
    \mathtt{proj}_i\,(t@\mathcal{U}^\gamma)\\
x@\mathcal{U}^\gamma &= \gamma\,x \\
(\lambda x.t)@\mathcal{U}^\gamma &=
    \mathtt{lam}\,\left(t@\mathcal{U}^{\left[
        \vec{y}\coloneqq
            \mathtt{comp}\,
                \mathtt{proj}_1(\gamma\,\vec{y}),\;
        x\coloneqq\mathtt{proj}_2
    \right]}\right)\\
(t\,s)@\mathcal{U}^\gamma &= \mathtt{comp}\,
    \mathtt{eval}\,(\mathtt{pair}\,
        (t@\mathcal{U}^\gamma)\,
        (s@\mathcal{U}^\gamma))
\end{align*}
Where universe operations \texttt{cnst},\,\texttt{pair} and others
are taken from $\mathcal{U}$.
\end{definition}

\begin{definition}
    Relation $\to_R$ is STLC-congruent if the following rules are
    admissible:
    \begin{mathpar}
        \inferrule{s\to_R s'}{(s,t)\to_R(s',t)}
        \and\inferrule{t\to_R t'}{(s,t)\to_R(s,t')}
        \and\inferrule{t\to_R t'}{\pi_i t\to\pi_i t'}\\
        \and\inferrule{t\to_R t'}
            {\lambda x.t\to_R\lambda x.t'}
        \and\inferrule{s\to_R s'}{s\,t\to_R s'\,t}
        \and\inferrule{t\to_R t'}{s\,t\to_R s\,t'}
    \end{mathpar}
\end{definition}

\begin{definition}
    Define \texttt{let}-reduction $(\to_\mathrm{let})$ to be the
    least STLC-congruent relation where rules $(\to_\mathtt{let})$
    from figure \ref{stlc-r-comp} are admissible.
\end{definition}

\begin{definition}
    Define $(\to)$ to be the STLC-congruent closure of
    \texttt{let}-reduction and rules $(\to_\beta)$ from figure
    \ref{stlc-r-comp}.
\end{definition}

\subsection{Metatheoretical properties of STLC-R}

Here we only study some basic properties of a system to make sure
that everything works as expected: computations don't diverge; types
are preserved by all computations; all of the types can be inferred
if you wish to do so; everything terminates and, modulo reduction of
\texttt{let}-bindings, this is the same old STLC we started with.
Proofs are mostly following standard techniques, since rules for
reinterpretation are structured to directly plug into proofs of
properties of STLC.

\begin{theorem}[Church-Rosser theorem for STLC-R]
    Reflexive transitive closure of $(\to)$ has Church-Rosser
    property.
\end{theorem}
\begin{proof}
    Classic proof via parallel reduction relation \cite{CRproof}.
    Both reductions outside and inside \texttt{let}-reduction have to
    be parallelized in order to prove the theorem.
\end{proof}

\begin{theorem}[Type preservation for STLC-R]
    The following is true:
    \begin{enumerate}
        \item If $\mathcal{U}=\langle\ldots\rangle$
            is a correct universe,
            $\gamma\,x:C[\Gamma_\mathcal{U},\Gamma(x)_\mathcal{U}]$
            for all $x\in\Gamma$ and $\Gamma\vdash t:T$, then
            $t@\mathcal{U}^\gamma:
                C[\Gamma_\mathcal{U},T_\mathcal{U}]$.

        \item If $s:\sigma$ and $\Gamma,x:\sigma\vdash t:T$, then
            $\Gamma\vdash t\triangleleft_x s:T$.

        \item If $\Gamma\vdash t:T$ and $t\to t'$, then
            $\Gamma\vdash t':T$.
    \end{enumerate}
\end{theorem}
\begin{proof}
    By structural induction on terms and $(\to)$. It is important
    that, due to the way computations are structured,
    $t\triangleleft_x s$ is always applied to closed $s$, therefore
    $t@\mathcal{U}$ is always applied to $t$ that is
    $\forall$-closed.
\end{proof}

\begin{theorem}[Strong normalization for STLC-R]
    Closed terms typable in STLC-R are strongly normalizing with
    respect to $(\to)$.
\end{theorem}
\begin{proof}
    By Tait's method of logical relations. Let $SN$ be a set of
    strongly normalizing terms. Let there be the following relation
    $\mathcal{R}\subseteq\mathsf{PolyType}\times\mathsf{Term}$:

    \begin{itemize}
        \item $t\in\mathcal{R}_\alpha,\mathcal{R}_\top$ iff
            $t\in SN$;
        \item $t\in\mathcal{R}_{T_1\times T_2}$ iff
            $\pi_i\,t\in\mathcal{R}_{T_i}$ for $i=1,2$;
        \item $t\in\mathcal{R}_{S\to T}$ iff for all
            $s\in\mathcal{R}_S$, $t\,s\in\mathcal{R}_T$;
        \item $t\in\mathcal{R}_{\forall\vec{\alpha}.T}$ iff
            $t\in SN_\mathtt{let}$ and for all
            $\mathcal{U}\in\mathcal{R}_{univ}$ and $\vec{S}$,
            $nf_\mathtt{let}(t)\,@\,\mathcal{U}\in\mathcal{R}_{
                \mathcal{U}\llbracket T\rrbracket[
                    \vec{\alpha}\coloneqq\vec{S}]}$.
    \end{itemize}

    $\mathcal{U}\in\mathcal{R}_{univ}$ is defined similarly by
    requiring that all terms inside $\mathcal{U}$ belong to
    respective $\mathcal{R}$-s.

    Then we can prove our main lemma by mutual induction on typing
    rules:

    \begin{itemize}
        \item If $\Gamma\vdash t:T$ and $\gamma$ is a substitution
            s.t. $\gamma\,x\in\mathcal{R}_{\Gamma(x)}$ for all
            $x\in\Gamma$ then $\gamma\,t\in\mathcal{R}_T$.
        \item If $\forall\Gamma\vdash t:T$, $\vec{\alpha}=fv(T)$ and
            $\gamma$ is a substitution s.t.
            $\gamma\,x\in\mathcal{R}_{\Gamma(x)}$ for all
            $x\in\Gamma$ then
            $\gamma\,t\in\mathcal{R}_{\forall\vec{\alpha}.T}$.
    \end{itemize}

    Noticing $\mathcal{R}_T\subseteq SN$ completes the proof.
\end{proof}

\begin{theorem}[Conservativity over STLC]
    Normal forms of closed terms typable in STLC-R are typable in
    STLC.
\end{theorem}
\begin{proof}
    Thanks to type preservation, normal form $nf(t)$ is typable in
    STLC-R. Normal form cannot contain let-bindings by definition of
    $(\to)$, therefore each subterm is $\forall$-closed. Therefore
    $nf(t)$ has no reinterpretation subterms.
\end{proof}

\begin{theorem}
    Global type inference for STLC-R is decidable.
\end{theorem}
\begin{proof}
    It is enough to extend any unification-based type inference
    algorithm for Hindley-Milner type system to support universes and
    reinterpretation.
    \begin{itemize}
        \item Since universes are provided explicitly, it is enough
            to \textit{check} them.
        \item Since polymorphic functions do not depend on local
            context, their type is fully resolved when it's added to
            context.
        \item Once universe is checked in $x\,@\,\mathcal{U}$, we can
            directly compute $\mathcal{U}\llbracket T\rrbracket$
            (since $\mathcal{U}$ is known and $T$ is fully resolved)
            and instantiate it with fresh metavariables in $\vec{S}$.
    \end{itemize}
\end{proof}

\subsection{Useful reinterpretations inside STLC-R}
\label{STLC-R-applications}

The idea to interpret STLC terms in cartesian closed categories is
really not new and is already thoroughly explored by Conal Elliott in
``Compiling to Categories'' \cite{Elliott:Compcat2017}; our addition
here is that we internalized reinterpretation inside the calculus
without working through compiler plugin. So all of the main
applications from the cited paper are still applicable here whenever
they are expressible as universes in STLC-R. As STLC-R is not the
main focus of the paper, we will not be reiterating the examples
here, but we would like to add just one more example.

Following Moggi \cite{MOGGI199155}, we can use monads to model side
effects in a language. What is new here is that instead of
parameterizing a function by arbitrary monad, we can simply
\textit{reinterpret} pure terms as if they were run in monadic
context. Therefore, our universe is a Kleisli category; however, it
is only cartesian, but in general not cartesian \textit{closed}. The
problem is the following operation (here $F$ is a monad in question,
$E$ is whatever exponential object we were to use):

\[
    \mathtt{lam}:(\alpha\times\beta\to F[\gamma])
        \to(\alpha\to F[E[\gamma,\beta]])
\]

In other words, $F$ has to \textit{distribute} over an arrow. Well,
at least arrow type distributes over itself:

\begin{align*}
    \mathcal{U}_E &= \langle\\
        &\Lambda01.\,0\to E\to1,\\
        &\lambda x e.x,\\
        &\lambda gfxe.g(fxe)e,\\
        &\top,\\
        &\lambda e.(),\\
        &\Lambda01.\,0\times1,\\
        &\lambda fgxe.(fxe,gxe),\\
        &\lambda pe.\pi_1\,p,\\
        &\lambda pe.\pi_2\,p,\\
        &\Lambda01.\,0\to1,\\
        &\lambda fxey.f(x,y)e,\\
        &\lambda pe.(\pi_1\,p)\,(\pi_2\,p)\\
    \rangle
\end{align*}

Interestingly enough, if, using Curry-Howard correspondence, we take
closed STLC term $t$ to represent proof of an intuitionistic
proposition, $t\,@\,\mathcal{U}_E$ corresponds to hypothesis
elimination in intuitionistic logic.

We \textit{could} write down a universe for arbitrary monad if we had
a capability to interpret terms which do not use abstractions in
arbitrary cartesian category, which is actually achievable with
\cite{Elliott:Compcat2017}. We leave study of such ``partial''
interpretation for future work.

\section{Dependently typed reinterpretable language: MLTT-R}
\label{MLTT-R}

Now we are ready to take on the main challenge: design a minimal
dependently typed reinterpretable language, and provide internal
models of itself to yield universe polymorphism, external
parametricity and do synthetic mathematics inside it. As a starting
point, we take Martin-L\"of type theory with $\Pi-$,$\Sigma-$,
$\mathrm{Id}-$ types and infinite cumulative hierarchy of universes.
The problems are similar to the ones we had with simply-typed
calculus:

\begin{itemize}
    \item How to reinterpret free variables?
    \item How to type universes?
    \item Now, there is a way to enforce that internal model
        satisfies coherence conditions on their operations. How to
        structure computation in such a way that it doesn't break
        Church-Rosser property?
    \item How to respect conversion rule?
\end{itemize}

As before, we sidestep the question with free variables by
considering a term for reinterpretation iff it is closed in local
context; thus reinterpretable terms form a separate sort which we
denote by $\square$ since it makes it really close to modal type
theories. To type universes, we have to consider:

\begin{itemize}
    \item How to provide a ``universe'' for each level in the
        interpreted term?
    \item Typing rules for MLTT are necessarily mutually recursive
        with rules of context well-formedness. Which framework to use
        to capture this in a model?
\end{itemize}

To deal with both problems, we use Categories-with-Families
\cite{castellan2020cwf} as one of the most convenient models of
dependent type theories in use today.

The problem with Church-Rosser is solved in absolutely the same way
as with simply-typed case: we just do not perform any
$\beta$-reductions inside a let-binding, only inline \texttt{let}
declarations before performing a substitution.

The problem with conversion rule is the following: in definition of a
universe, we can now ask for propositional proofs of coherence. Thus
types which were \textit{definitionally} equal before
reinterpretation are now only \textit{propositionally} equal. To
solve this issue, we start by first pinpointing the locations where
conversion rule is actually used; in order to do this, we
\textit{need} to use bidirectional type checking. After translation,
point of conversion is wrapped in a \texttt{rewrite} call.

But how do we synthesize a propositional proof of equivalence? Here
we translate a chain of reductions performed during checking of
definitional equality into a chain of propositional equivalences
taken from coherence conditions proven in the definition of a model.

\begin{figure}
    \centering
\begin{bnf}(comment={//}, single-line-delim={\&})
    $i$ // \textsf{Level} ::=
    | $l$ // variable
    | $0$ // ground level
    | $suc\,i$ // next universe level
    ;;
    $S,T,s,f$ // \textsf{Term} ::=
    | $U_i$ // universe
    | $x$ // variable
    | $\Sigma_{x:S}T(x)$ // dependent sum type
    | $(s, t)$ // dependent sum introduction
    | $\pi_i t$ // projection out of dependent sum
    | $\Pi_{x:S}T(x)$ // dependent product type
    | $\lambda x.t$ // abstraction
    | $t\,s$ // application
    | $s=_T t$ // propositional equality type
    | $\mathtt{refl}\,t$ // reflexivity proof
    | $\mathtt{rewrite}_{x\,y.T(x,y)}s\,t$ // rewrite
    | $\mathtt{let}\,x=s\,\mathtt{in}\,t$
        // reinterpretable definition
    | $x\,@\,\mathcal{U}$ // reinterpretation
    ;;
    $\mathcal{U}$ // \textsf{Universe} ::=
    | $\langle \mathtt{C},\ldots \rangle$
        // definition of a model (see below)
\end{bnf}
    \caption{MLTT-R grammar}
    \Description{MLTT-R grammar}
    \label{mltt-r-grammar}
\end{figure}

\begin{figure}
    \centering
\begin{mathpar}
    \inferrule{}{\cdot\,\mathrm{ctx}}
    \and\inferrule{
        \Gamma\,\mathrm{ctx}\\
        \Gamma\vdash T\Rightarrow U_i
    }{
        \Gamma,x:T\,\mathrm{ctx}\\
        \Gamma,x:\square T\,\mathrm{ctx}
    }
    \and\inferrule{}{
        \Gamma\vdash U_i\Rightarrow U_{\mathrm{suc}\,i}
    }
    \and\inferrule
        {\Gamma\vdash t\Rightarrow U_i\\ i\le j}
        {\Gamma\vdash t\Leftarrow U_j}
    \and\inferrule{}
        {\Gamma,x:T,\Delta\vdash x\Rightarrow T}
    \and\inferrule{
        \Gamma\vdash S\Leftarrow U_i\\
        \Gamma,x:(S:U_i)\vdash T\Leftarrow U_i
    }{\Gamma\vdash\Sigma_{x:S}T(x)\Leftarrow U_i}
    \and\inferrule{
        \Gamma\vdash s\Leftarrow S\\
        \Gamma\vdash t\Leftarrow T(s)
    }{\Gamma\vdash(s, t)\Leftarrow\Sigma_{x:S}T(x)}
    \and\inferrule
        {\Gamma\vdash t\Rightarrow\Sigma_{x:S}T(x)}
        {\Gamma\vdash \pi_1 t\Rightarrow S}
    \and\inferrule
        {\Gamma\vdash t\Rightarrow\Sigma_{x:S}T(x)}
        {\Gamma\vdash \pi_2 t\Rightarrow T(\pi_1 t)}
    \and\inferrule{}{\pi_i((t_1,t_2):\Sigma_{x:S}T(x))\to_\beta t_i}
    \and\inferrule{
        \Gamma\vdash S\Leftarrow U_i\\
        \Gamma,x:(S:U_i)\vdash T\Leftarrow U_i
    }{\Gamma\vdash\Pi_{x:S}T(x)\Leftarrow U_i}
    \and\inferrule
        {\Gamma\vdash S\Rightarrow U_i\\
        \Gamma,x:S\vdash t\Leftarrow T(x)}
        {\Gamma\vdash\lambda x.t\Leftarrow
            \Pi_{x:S}.T(x)}
    \and\inferrule{
        \Gamma\vdash t\Rightarrow\Pi_{x:S}T(x)\\
        \Gamma\vdash s\Leftarrow S
    }{\Gamma\vdash t\,s\Rightarrow T(s)}
    \and\inferrule{}{
        (\lambda x. t : \Pi_{x:S}T(x))s\to_\beta
        t[x\coloneqq(s:S)]
    }
    \and\inferrule{
        \Gamma\vdash T: U_i\\
        \Gamma\vdash s,t\Leftarrow T
    }{\Gamma\vdash s=_T t: U_i}
    \and\inferrule
        {\Gamma\vdash t\Rightarrow T}
        {\Gamma\vdash\mathtt{refl}\,t\Rightarrow t=_T t}
    \and\inferrule{
        \Gamma\vdash s\Rightarrow s_1=_S s_2\\
        \Gamma,x:S,y:s_1=_S x\vdash
            T(x,y)\Rightarrow U_i\\
        \Gamma\vdash t\Leftarrow
            T(s_1, \mathtt{refl} (s_1 : S))
    }{\Gamma\vdash\mathtt{rewrite}_{x\,y.T(x,y)}s\,t
        \Rightarrow T(s_2, s)
    }
    \and\inferrule{}{
        \mathrm{rewrite}_{x\,y.T(x,y)}
            \,(\mathrm{refl} \,(s : S))\,t\to_\beta t
    }
    \and\inferrule{
        \square\Gamma\vdash s\Rightarrow S\\
        \square\Gamma,x:\square S,\Delta\vdash t:T
    }{
        \square\Gamma,\Delta\vdash
        \mathtt{let}\,x=s\,\mathtt{in}\,t:T
    }
    \and\inferrule
        {s\to_\mathtt{let} s'}
        {\mathtt{let}\,x=s\,\mathtt{in}\,t
        \to_\mathtt{let}
        \mathtt{let}\,x=s'\,\mathtt{in}\,t}
    \and\inferrule
        {\mathtt{let}\not\in s}
        {\mathtt{let}\,x=s\,\mathtt{in}\,t
        \to_\mathtt{let}t[x\coloneqq s]}
    \and\inferrule{
        \square\Gamma\vdash x\Rightarrow\square T\\
        \mathcal{U}\,\mathrm{model}
    }{x\,@\,\mathcal{U}\Rightarrow\mathcal{U}\llbracket T\rrbracket}
    \and\inferrule{
        \mathtt{C}\Rightarrow U_0\\
        \mathtt{Sub}\Rightarrow C\to C\to U_0\\
        \mathtt{*}\Rightarrow C\\
        \mathtt{Ty}\Rightarrow C\to U_0\\
        \ldots
    }{\langle \mathtt{C},\ldots \rangle\,\mathrm{model}}
\end{mathpar}
    \caption{MLTT-R typing rules}
    \Description{MLTT-R typing rules}
    \label{mltt-r-typing}
\end{figure}

To actually compute interpretation in MLTT-R, we need two mutually
recursive functions: $t\,@^\gamma_{\Leftarrow T}\mathcal{U}$ and
$t\,@^\gamma_{\Rightarrow}\mathcal{U}$ for checkable and synthesizing
terms, respectively. \textbf{We're sorry to be glossing over details here and will provide all definitions in the updated version of the paper.}

\subsection{Metatheoretical properties of MLTT-R}

Once again, we only try to make sure of the most essential
properties: that computations do not diverge, that types during
computations are preserved, all computations terminate and that
resulting type system is, in essence, the same as base MLTT modulo
inlining of \texttt{let}-bindings.

\begin{theorem}[Church-Rosser theorem for STLC-R]
    Reflexive transitive closure of $(\to)$ has Church-Rosser
    property.
\end{theorem}
\begin{proof}
    Classic proof via parallel reduction relation \cite{CRproof}.
    Both reductions outside and inside \texttt{let}-reduction have to
    be parallelized in order to prove the theorem.
\end{proof}

\begin{theorem}[Type preservation for MLTT-R]
\end{theorem}
\begin{proof}
    By structural induction on terms and $(\to)$. It is important
    that, due to the way computations are structured,
    $t\triangleleft_x s$ is always applied to closed $s$, therefore
    $t@\mathcal{U}$ is always applied to $t$ that is
    $\square$-closed.
\end{proof}

\begin{theorem}[Strong normalization for MLTT-R]
    Closed terms typable in MLTT-R are strongly normalizing with
    respect to $(\to)$.
\end{theorem}
\begin{proof}
    Tait's method of logical relations for MLTT with modifications
    similar to ones performed in a proof of strong normalization for
    STLC-R should suffice. However, we're not completely sure.
\end{proof}

\begin{theorem}[Conservativity over MLTT]
    Normal forms of closed terms typable in MLTT-R are typable in
    MLTT.
\end{theorem}
\begin{proof}
    Thanks to type preservation, normal form $nf(t)$ is typable in
    MLTT-R. Normal form cannot contain let-bindings by definition of
    $(\to)$, therefore each subterm has no $\square$ in context.
    Therefore $nf(t)$ has no reinterpretation subterms.
\end{proof}

\subsection{Useful reinterpretations inside MLTT-R}
\label{MLTT-R-applications}

Here we list some of the language features which reinterpretation
enables without any additional language support.

\subsubsection{Crude, but effective stratification}
\label{MLTT-R-displacement}

Due to McBride \cite{McBride:Crude2002}, a technique of crude, but
effective stratification works as follows: instead of providing
complete support for level polymorphism, let's consider a
\textit{displacement operator} which, given a term $t$ using
universes $U_i,U_j,U_k,\ldots$, produces a term which is equal to $t$
except for universes which are replaced with
$U_{i+1},U_{j+1},U_{k+1},\ldots$. The displaced term is also
well-typed, however it is defined <<one level higher>>. Notice that
displacement operator is \textit{definable} in MLTT-R as
$\_\,@\,\mathcal{U}_\mathrm{McBride}$ where

\[\mathcal{U}_\mathrm{McBride} = \langle
  \Lambda l.U_{suc\,l}, \ldots
\rangle\]

\subsubsection{Internal external parametricity}
\label{MLTT-R-parametricity}

Demonstrated by Reynolds \cite{reynolds1983types}, parametricity is a
property of a type system that, in simple words, <<polymorphic
function cannot meaningfully depend on generic inputs>>. Thus any
polymorphic function of type $\sigma$ satisfies a certain <<free
theorem>> $\llbracket \sigma \rrbracket$ defined via relational
interpretation of a type.

In recent works \cite{Bernardy:Parametricity2010}
\cite{Altenkirch:Parametricity2024}, a streamlined definition is
given for parametricity interpretation of dependently-typed systems
dubbed <<external parametricity>>. Once again, this relational
interpretation is definable in our system \textit{internally} with
the help of a relational model $\mathcal{U}_\mathrm{para}$. Since
interpreted functions do not depend on local context, we sidestep
issues with relational context highlighted by Altenkirch and others
\cite{Altenkirch:Parametricity2024}, thus result of the
interpretation is a proper proof of a free theorem.

\subsubsection{Synthetic mathematics} \label{MLTT-R-synthetic}

As thoroughly discussed in a recent paper by Nawrocki and others
\cite{Nawrocki:Lean2026}, thanks to MLTT being an internal language
of certain $(\infty,1)$-categories, we can do \textit{synthetic
mathematics} inside it: prove a theorem inside MLTT about types;
interpret it in a model of MLTT and yield a proof about objects in
that model. It is easy to see that reinterpretation operator we have
in MLTT-R is sufficient to do synthetic mathematics in it and
immediately apply results interpreted in a suitable model.

\section{Related work}

To the best of our knowledge, we didn't find any research which
explores the topic at hand, but the idea is certainly in the air, as
illustrated by the following brilliant papers.

Original work by Lambek \cite{Lambek1980} sets the tone of the future
research, establishing correspondence between type systems based on
lambda calculus with various kinds of categories, starting with
correspondence between simply typed lambda calculus and cartesian
closed categories.

Conal Elliott's ``Compiling to Categories''
\cite{Elliott:Compcat2017} shows how to interpret (some subset of)
Haskell in arbitrary \texttt{Category} (again, definable in Haskell)
via GHC compiler plugin. That paper is a direct inspiration for this
one. There is an important difference, though: reinterpretation is
internal to our type systems, unlocking new patterns in the usage of
reinterpretation, including (but not limited to) universe
polymorphism, external parametricity and more. However, Conal's
approach is more flexible since his translation to categories
requires desired category to only implement primitives actually used
for translation of a particular term at hand.

In a recent paper ``A Certifying Proof Assistant for Synthetic
Mathematics in Lean'', Wojciech Nawrocki and others
\cite{Nawrocki:Lean2026} adapt Conal's approach to Lean's system with
dependent types, using Categories-with-Families to represent internal
models and employing Lean's powerful metaprogramming facilities
instead of a compiler plugin. In doing so, they are able to provide a
framework for doing synthetic mathematics in Lean, extracting proofs
done in MLTT applied to more specific categories. Thus our results in
section \ref{MLTT-R-synthetic} are close to the results shown in
their paper. To add to that, their work is already implemented in
practical proof assistant. Yet, our type system involves no
metaprogramming since reinterpretation is internal to it, and it
handles conversion rule gracefully. In addition, the focus is a bit
different, where we use reinterpretation to express other polymorphic
features of a language.

Conor McBride's ``Crude but effective stratification''
\cite{McBride:Crude2002} (systemic treatment of this idea is given by
Favonia and others \cite{Favonia:Displacement2023}) introduces an
eye-opening idea of using \textit{displacement} as a lightweight form
of universe polymorphism: to reuse term typed at a lower universe
level, lift all levels in a term by adding a constant shift $s$. From
the point of view advocated by this paper, you can see it as an act
of reinterpreting an expression in an internal model of type system
where universe at level $i$ is interpreted as a universe at level
$i+s$ (section \ref{MLTT-R-displacement}).

In ``Parametricity and Dependent Types'' by Jean-Philippe Bernardy,
Patrik Jansson and Ross Paterson \cite{Bernardy:Parametricity2010},
authors introduce a direct translation procedure to yield
parametricity theorems and proofs from polymorphic definitions in a
dependently typed system. The definition of <<external
parametricity>> we follow, though, is given in ``Internal
Parametricity, without an Interval'' by Thorsten Altenkirch, Yorgo
Chamoun, Ambrus Kaposi and Michael Shulman
\cite{Altenkirch:Parametricity2024}. Once again, we can reformulate
the translation they give as reinterpretation in an internal model.
Altenkirch and others noted a potential problem with this approach:
if we apply the external translation to the judgment with nonempty
context, context becomes parametric, too; that's the reason they add
\texttt{rel} operation and yield type system with
\textit{computational properties of higher observational type
theory}. We sidestep the issue with context since reinterpretation
is applied only to terms whose context is also <<reinterpretable>>
and is not affected by the translation
(section \ref{MLTT-R-parametricity}).

\section{Conclusions and further work}

In retrospect, we managed to define two self-interpreting type
systems, STLC-R and MLTT-R, meaning that they have a notion of their
own <<internal models>> which they then can use to reinterpret their
own terms inside such models. In doing so, such type systems get new
features which previously were defined as separate builtin quirks of
a system. Basic metatheoretical properties are proven; for future
work we have the following plans:

\begin{itemize}
    \item Provide mechanized proofs of metatheoretical properties of
        systems;
    \item Provide implementation of MLTT-R in a practical proof
        assistant;
    \item Handle <<partial>> interpretation \`a-la Conal Elliott's
        ``Compiling to Categories'' \cite{Elliott:Compcat2017};
    \item Make equality in MLTT-R more extensional, most probably by
        making type system closer to Homotopy type theory
        \cite{hottbook};
    \item Relax context constraints in reinterpretable terms.
\end{itemize}

\bibliographystyle{ACM-Reference-Format}
\bibliography{imtt}

\end{document}
